\section{Discussion}\label{sec:discussion}

\subsection{Verdicts on the hypotheses}\label{sec:mapping}

The evidence supports H1, H2, H3 and H4a with conventional precision and H4b and H4c in direction only; the efficient-reading rival of H3 is rejected (Table~\ref{tab:hyp}). Two readings need qualification. The block predicts the signal and, once the signal is in the model, adds nothing to the prediction of the disposal or its timing, as it would if the stake works by shaping the intention the filing discloses; it remains an association, since we have no instrument for the stake. The reform results rest on a difference against an untreated form; the three measures arrived within one quarter, so the difference-in-differences is the package's, and the appraisal-right and protection-block tests use \AppraisalPrices{} and \ProtPostN{} post-reform filings and a short disposal record, the appraisal slope reaching conventional significance only in the 21-day window. They establish the direction of each channel, not its size in steady state. Beside the independently built sample of \citeA{kang2025} (Supporting Information, Table S5) the two papers agree on every fact both measure---a small positive announcement return, a negative return over the following year concentrated in disposed subsidiaries and small-stake parents, a disposal rate near seventy percent, a listing rate near four percent---and differ in what they add: their contribution is to locate the loss at the disposal, ours to show that the disposal could be forecast from the controller's own filing at the time the minority voted on it.

\subsection{Theoretical implications}\label{sec:theoretical}

Three of the four strands of Section~\ref{sec:theory} gain something from the pattern, and the fourth loses a presumption. For the principal--principal literature, which has treated the controlling owner's intention as unobservable and inferred it from the expropriating outcome \cite{bertrand2002, bae2002, baek2006, buchuk2014, chen2024}, the result is that for one class of transactions the outcome was forecastable before it occurred from the document the controller was legally required to produce to obtain the minority's vote, and that the forecast moved with the controller's stake as \eqref{eq:pp} implies. The disclosure strand explains why the forecast was there: a disposal-oriented division has to be described to the counterparties it needs, and counterparty-addressed statements---a named buyer, a stated plan to attract outside capital, a charter clause that binds---carry the information that the form's purpose categories do not \cite{crawford1982, farrell1996, jensen1976}; the nested models put the two apart in one exhibit, the categories lifting the AUC from \NestedMZero{} to \NestedMOne{} and the reading from \NestedMTwo{} to \NestedMThree{}. The disclosure literature's question about expropriation has usually been whether a mandate changes conduct \cite{benshahar2014}; the prior question is whether the disclosure already produced contains the forecast, and whether it is priced. Here the filings contained it, and the announcement returns show no sign that it was priced. That is the loss for the fourth strand, the efficient-markets presumption that attaches to event studies of divestitures \cite{schipper1986, nanda1991, slovin1995, hulburt2002}: in a market in which the disposal is at the controller's discretion, a positive announcement return is consistent with investors pricing the stated growth story and not the disposal the filing's details implied. Investors' understanding is not observed. The data show that a reader with no stake and no knowledge of the later history could extract the forecast from the filing, and that prices moved with the disposal rather than with the filing.

\subsection{Implications for investors, boards and controlling owners}\label{sec:practitioner}

For an institutional investor a \captive{} announcement is a dated event with a measurable exposure: moving from the retention to the disposal end of the signal places the parent \RankSpecThreePctl{} percentiles lower in its year's return ranking, the exposure is concentrated in the first year and in parents with blocks below about \LowBlockMax{} percent, and within a year \KMHighOneY{} percent of the subsidiaries in the top tercile of the signal have changed hands. The division filing is the one document, and the general meeting the one occasion, on which the disposal can be contested, and the question to put to it is whether the new company will remain wholly owned; the sentences the models quote---Doosan's statement that the parent would sell its entire stake to an affiliate, CJ ENM's ``attracting outside investment \ldots{} and the sale of shares where necessary''---are of the kind an analyst can find in an afternoon, and a stewardship policy can require that finding before the vote. The same reading is available to a proxy adviser or to the exchange's listing committee, and the appraisal price, which is in the filing, tells a dissenting shareholder on the day how far the price of dissent stands above or below the market.

For boards and controlling owners the finding cuts the other way. A parent that intends to keep the subsidiary pays nothing to say so in a way that binds, and the market reads the difference: the protection block's content moves with the signal and with the labelled outcomes, a charter clause or a stated horizon lowering the signal a reader extracts where a recital of the form does not. A board that means to retain can therefore lower the disposal risk its own filing discloses, and the one-year drift that accompanies it, by writing a commitment rather than a purpose; a board that means to dispose cannot, and after 2022 writes its intention into a document the exchange will read if the subsidiary applies to list within five years. For the controlling owner the result quantifies the private cost of the stake: a block ten percentage points smaller goes with a signal \ScoreOnBlockTenPPSDAbs{} standard deviations higher and, across terciles, with \LowBlockDisp{} rather than \HighBlockDisp{} percent of subsidiaries disposed of, which is what \eqref{eq:pp} predicts and what a minority investor will now price into the parent of a low-block controller who divides.

\subsection{Policy implications}\label{sec:policy}

For regulators, each channel is judged by what it changed. The information channel moved the filings and their reception---relative to equity spin-offs a \captive{} announcement was \DiDCxP{} percentage points better received after the package---while the mean forecast moved from \ScorePre{} to \ScorePost{} and the early-disposal rate from \EarlyPre{} to \EarlyPost{} percent, changes inside what \ProtPostN{} filings can detect. The compensation channel is small: the appraisal right is near the money at the announcement, its moneyness carries the sign of a put but is priced imprecisely, and a dissenting shareholder who exercised it would have beaten holding the parent in \ExerciseBeats{} percent of cases by a small margin. The deterrence and bonding channel is the one that reaches the controller's intention, and it does so through commitments that bind rather than statements of intent. A design implication follows from the timing of the loss, untested here. The right to dissent and the price of dissent are fixed at the division, when the market's reaction is positive, and the loss arrives at the disposal, when no right exists. A right exercisable at the disposal, or a vote of the parent's shareholders on a disposal within some years of a \captive{}, would act at the step where the loss occurs, and the forecastability we document implies that such a vote would be informed; it would also give dissenters a hold-up over a transaction the board may need to close, a cost our data do not measure. Where disclosure remains the instrument, the statements can be audited against conduct: a protection statement in a filing whose subsidiary applies to list within the review period can be compared with what the parent did, and the comparison published.
